\section{Schanuel Lemma}

\begin{theorem}[Schanuel Lemma]
    Let
    \[
        0\longrightarrow K \longrightarrow P \xrightarrow{\,\lambda\,} M \longrightarrow 0
    \]
    and
    \[
        0\longrightarrow K' \longrightarrow P' \xrightarrow{\,\lambda'\,} M \longrightarrow 0
    \]
    be short exact sequences of left $R$-modules, with $P$ and $P'$ projective. Then
    \[
        K'\oplus P \cong K\oplus P'.
    \]
\end{theorem}

\begin{proof}
    Define
    \[
        Y:=P\times_M P'
        =\{(p,p')\in P\oplus P' \mid \lambda(p)=\lambda'(p')\}.
    \]
    Let
    \[
        \mu:Y\to P,\qquad \mu(p,p')=p,
    \]
    and
    \[
        \mu':Y\to P',\qquad \mu'(p,p')=p'.
    \]
    Then $Y$ fits into the pullback diagram
    \[
        \begin{tikzcd}
            Y \arrow[r, "\mu'"] \arrow[d, "\mu"']
            & P' \arrow[d, "\lambda'"] \\
            P \arrow[r, "\lambda"']
            & M
        \end{tikzcd}
    \]
    that is, the square commutes and $Y$ is the fiber product of $\lambda$ and $\lambda'$.

    We claim that there are short exact sequences
    \[
        0\longrightarrow K' \xrightarrow{\,i'\,} Y \xrightarrow{\,\mu\,} P \longrightarrow 0,
        \qquad i'(k')=(0,k'),
    \]
    and
    \[
        0\longrightarrow K \xrightarrow{\,i\,} Y \xrightarrow{\,\mu'\,} P' \longrightarrow 0,
        \qquad i(k)=(k,0).
    \]

    First,$\mu$ is surjective. Indeed, let $p\in P$. Since $\lambda'$ is surjective,
    there exists $p'\in P'$ such that
    \[
        \lambda'(p')=\lambda(p).
    \]
    Then $(p,p')\in Y$, so $\mu(p,p')=p$.

    Next,
    \[
        \ker(\mu)
        =\{(0,p')\in Y \mid \lambda'(p')=0\}
        =\{(0,k') \mid k'\in K'\}
        =\operatorname{Im}(i').
    \]
    Hence
    \[
        0\longrightarrow K' \xrightarrow{\,i'\,} Y \xrightarrow{\,\mu\,} P \longrightarrow 0
    \]
    is exact.

    Similarly,$\mu'$ is surjective, and
    \[
        \ker(\mu')
        =\{(p,0)\in Y \mid \lambda(p)=0\}
        =\{(k,0) \mid k\in K\}
        =\operatorname{Im}(i).
    \]
    Thus
    \[
        0\longrightarrow K \xrightarrow{\,i\,} Y \xrightarrow{\,\mu'\,} P' \longrightarrow 0
    \]
    is exact.

    Since $P$ is projective and $\mu:Y\to P$ is an epimorphism, there exists
    an $R$-homomorphism $s:P\to Y$ such that
    \[
        \mu\circ s=\mathrm{id}_P.
    \]
    This is expressed by the commutative diagram
    \[
        \begin{tikzcd}
            & Y \arrow[d, two heads, "\mu"] \\
            P \arrow[ur, dashed, "s"] \arrow[r, "\mathrm{id}_P"']
            & P
        \end{tikzcd}
    \]
    so the sequence
    \[
        0\longrightarrow K' \longrightarrow Y \xrightarrow{\,\mu\,} P \longrightarrow 0
    \]
    splits. Therefore,
    \[
        Y\cong K'\oplus P.
    \]

    Likewise, since $P'$ is projective and $\mu':Y\to P'$ is an epimorphism, there exists $s':P'\to Y$ such that
    \[
        \mu'\circ s'=\mathrm{id}_{P'},
    \]
    and hence
    \[
        0\longrightarrow K \longrightarrow Y \xrightarrow{\,\mu'\,} P' \longrightarrow 0
    \]
    splits. Therefore,
    \[
        Y\cong K\oplus P'.
    \]

    Combining the two decompositions of $Y$, we obtain
    \[
        K'\oplus P \cong Y \cong K\oplus P',
    \]
    and hence
    \[
        K'\oplus P \cong K\oplus P'.
    \]
\end{proof}



\begin{lemma}
    Let $M$ be a left $R$-module. Then
    \[
        \Hom_R({}_R R,M)\cong M
    \]
    as left $R$-modules.
\end{lemma}

\begin{proof}
    Define
    \[
        \Phi:\Hom_R({}_R R,M)\to M,\qquad f\mapsto f(1).
    \]
    This map is $R$-linear. Define also
    \[
        \Psi:M\to \Hom_R({}_R R,M),\qquad m\mapsto (r\mapsto rm).
    \]
    For $f\in \Hom_R({}_R R,M)$ and $r\in R$,
    \[
        \Psi(\Phi(f))(r)=r f(1)=f(r\cdot 1)=f(r),
    \]
    so $\Psi\Phi=\operatorname{id}$. Also
    \[
        \Phi(\Psi(m))=\Psi(m)(1)=1m=m,
    \]
    hence $\Phi$ and $\Psi$ are inverse isomorphisms.
\end{proof}

\begin{corollary}
    For any ring $R$ with $1$,
    \[
        \operatorname{End}_R({}_R R)\cong R^{\operatorname{op}}
    \]
    as rings.
\end{corollary}

\begin{proof}
    By the lemma, every $f\in \operatorname{End}_R({}_R R)$ is determined by $f(1)$, and
    \[
        f(r)=rf(1)
    \]
    for all $r\in R$. Thus
    \[
        \Theta:\operatorname{End}_R({}_R R)\to R^{\operatorname{op}},\qquad f\mapsto f(1)
    \]
    is bijective. Moreover, for $f,g\in \operatorname{End}_R({}_R R)$,
    \[
        (f\circ g)(1)=f(g(1))=g(1)f(1),
    \]
    which is multiplication in the opposite ring. Therefore $\Theta$ is a ring isomorphism.
\end{proof}

